\section{Results}
\label{sec:results}

\subsection{Roadmap of empirical tests}
\label{sec:roadmap}

The DSC framework predicts that a $1/\ln^2(t/t_P)$ relaxation should manifest at three independent observational scales: the fine-structure constant $\alpha$, the Hubble parameter $H_0$, and the CMB angular spectrum.
We organize the empirical content of this section into three blocks corresponding to those scales, with each block reporting both anchor-scale tests (where the doubly-logarithmic ansatz has its full dynamic range; Pillar~M, \cref{sec:pillarM}) and dense-data tests (where the ansatz is mathematically degenerate with constant; Pillar~M and \cref{sec:pillarN}).
The first two blocks consist entirely of analyses of published peer-reviewed observational data; the third block reports forward and inverse computations of the CMB lattice surrogate.

\begin{itemize}[leftmargin=2em]
  \item \textbf{Block 1 (\cref{sec:block_alpha}).} Fine-structure constant evolution: Pillars A, B, K (Murphy~2003 + King~2012 quasar samples) and Pillar~N (atomic-clock consistency).
  \item \textbf{Block 2 (\cref{sec:block_H0}).} Hubble parameter relaxation: Pillars C, D, E, F, G, H, I, J, M, O (anchor $H_0$ + cosmic chronometers + Pantheon$+$ + DESI BAO + leverage diagnostic + JWST forecast).
  \item \textbf{Block 3 (\cref{sec:block_cmb}).} CMB lattice forward modeling: Gaussian one-point statistics, ablation against dissipative baseline, lattice acoustic-peak structure, $\chi^2$ benchmark vs.\ CAMB Planck-2018, and differentiable inverse reconstruction with held-out hemisphere diagnostic.
\end{itemize}

The framework's strongest empirical case rests on Block 1 (Pillar~A: $5.6\sigma$ Keck-only quasar drift; Pillar~K: $+8.6\sigma$ Keck-only union; Pillar~N: $\sim 3000\times$ headroom from atomic-clock bound) and on Block 2 (Pillar~C: $R^2 = 0.91$ four-anchor $H_0$ relaxation; Pillar~M: $8.4\times$ leverage explanation of why dense Block-2 data is consistent with the anchor result; Pillar~O: $\gtrsim 11\sigma$ falsification forecast within $\sim 5$ years).
Block 3, in contrast, is a methodological proof of concept: the lattice forward model exhibits CMB-like phenomenology but does not reproduce the Planck $C_\ell$ at the level of individual multipoles (\cref{sec:camb}); the forward result is to be read in conjunction with the held-out hemisphere diagnostic (\cref{sec:inverse}) rather than as a quantitative match.

\subsection{Block 1: Fine-structure constant evolution}
\label{sec:block_alpha}

This block reports four tests of the predicted $1/\ln^2(t)$ drift in the fine-structure constant $\alpha$, using two independent peer-reviewed quasar absorption-line catalogs (Murphy, Webb \& Flambaum~2003 Keck/HIRES sample and King \emph{et al.}~2012 Keck+VLT sample), and one external constraint (Rosenband \emph{et al.}~2008 atomic clocks).
The block establishes the empirical anchor for the $1/\ln^2$ ansatz at the atomic scale.

\paragraph{Pillar A: Murphy 2003 128-absorber quasar drift, two error budgets (\texttt{notebooks/15, 25}).}
\label{sec:pillarA}
On the 128-absorber Keck/HIRES catalog of Murphy, Webb \& Flambaum~2003~\citep{Murphy2003} (the actual sample underlying Pillar~A in the companion preprint Zenodo \texttt{19218674}), we fit Wang's $1/\ln^2(t)$ drift model and the constant ($\Delta\alpha/\alpha = 0$) null.
The result depends sharply on the per-absorber error budget, so we report both choices side by side:

\begin{itemize}[leftmargin=2em]
  \item \textbf{Budget (i): per-absorber fit error only ($\sigma_{\rm fit}$, no $\sigma_{\rm rand}$).}
    This is the convention used in the companion preprint Zenodo~\texttt{19218674} that originally reported Pillar~A.
    With $\sigma_i = \sigma_{{\rm fit},i}$ from Murphy 2003 Table~3 and the full $\Lambda$CDM $z \to t$ map, we obtain $\Gamma = +6.42 \pm 1.14$ ($5.62\sigma$, $\Delta\chi^2 = 31.54$, reduced $\chi^2 = 1.52$ for DSC vs.\ $1.75$ for the null).
    This is the headline ``Pillar~A'' figure.
  \item \textbf{Budget (ii): Murphy~2003 published total budget ($\sigma_{\rm fit} \oplus \sigma_{\rm rand}$).}
    Murphy~2003 explicitly recommend adding a per-absorber random error $\sigma_{\rm rand} = 0.91 \times 10^{-5}$ (high-contrast) or $2.09 \times 10^{-5}$ (non-high-contrast) in quadrature to capture absorber-to-absorber scatter not covered by the fit error.
    With this published total budget the same data give $\Gamma = +5.80 \pm 2.08$ ($2.79\sigma$, $\Delta\chi^2 = 7.80$).
    The reduced $\chi^2$ falls to $0.71$ for DSC and $0.77$ for the null, indicating that the budget-(i) reduced $\chi^2$ of $1.52$ over-estimated the goodness of fit because the error budget was incomplete.
\end{itemize}

The two budgets agree on the central value of $\Gamma$ ($\sim 5$--$6$) but differ in the formal significance.
Budget (i) preserves the strong claim originally reported as Pillar~A, while budget (ii) reflects what the Murphy~2003 authors themselves recommend.
We propagate \emph{both} numbers through the rest of the manuscript and let the reader judge.
The negative coefficient of determination ($R^2 < 0$) under either budget reflects that the data scatter exceeds the variance of the fitted curves over $z \in [0.23, 3.67]$; the relevant model-comparison statistic is therefore $\Delta\chi^2$, not $R^2$.
\Cref{fig:alpha_quasar} shows the fit under budget (i).
Pillar~K below decomposes the apparent signal by spectrograph (Keck vs.\ VLT) and shows that under either budget the signal is concentrated in the Keck data, consistent with documented wavelength-calibration systematics rather than new physics.

\begin{figure}[htbp]
  \centering
  \includegraphics[width=\columnwidth]{fig18_alpha_drift_quasar.pdf}
  \caption{Quasar $\Delta\alpha/\alpha$ drift on the Murphy~2003 128-absorber Keck/HIRES sample under error budget~(i) (per-absorber fit error only).
  The DSC $1/\ln^2(t)$ model (red) achieves $\Delta\chi^2 = 31.54$ ($5.6\sigma$) over the constant null (gray); under Murphy's published $\sigma_{\rm rand}$-augmented budget~(ii) the same data give $\Delta\chi^2 = 7.80$ ($2.8\sigma$).
  Reproducible from \texttt{notebooks/15\_alpha\_drift\_quasar.ipynb} and \texttt{notebooks/25\_murphy03\_replication.ipynb}.}
  \label{fig:alpha_quasar}
\end{figure}
\paragraph{Pillar B: King~2012 295-absorber sample (\texttt{notebooks/17}).}
For consistency we re-ran the same comparison on the full published King \emph{et al.}~(2012) catalog (CDS J/MNRAS/422/3370 \texttt{tablea1}, $N = 295$ after standard outlier handling).
The DSC $p = 2$ model gives reduced $\chi^2 = 0.881$ vs $0.878$ for the constant null, $\Delta\chi^2 \approx 0$, and the BIC penalty for one extra parameter ($\Delta {\rm BIC} = +5.7$) decisively favors the null.
The joint quasar + atomic-clock~\citep{Rosenband2008} bound is $|\Gamma| < 0.30$, far below the $\Gamma \sim 6$ best-fit on the smaller subset (\cref{fig:king2012}).
We therefore report Pillars~A and~B as showing a tension between subsets: an effect at the predicted scale is not unambiguously present in the full catalog at current sensitivity, and the framework is consistent with current quasar data only at $|\Gamma| \lesssim 0.3$.

\begin{figure}[htbp]
  \centering
  \includegraphics[width=\columnwidth]{fig20_quasar_293_fit.pdf}
  \caption{Three-model comparison on the full King \emph{et al.}~(2012) 295-absorber sample.
  Top: data with best-fit DSC $p = 2$ curve and constant null.
  Bottom: $\chi^2(\Gamma)$ profile and joint clock + quasar bound $|\Gamma| < 0.30$.
  $\Delta {\rm BIC}({\rm DSC}, {\rm null}) = +5.7$ favors the null.
  Reproducible from \texttt{notebooks/17\_quasar\_293\_fit.ipynb}.}
  \label{fig:king2012}
\end{figure}
\paragraph{Pillar K: Resolving the Pillar~A vs Pillar~B tension with the Murphy~2003 sample.}
We were able to obtain, after the prior commit, the original 128-absorber catalog of Murphy, Webb \& Flambaum~2003~\citep{Murphy2003} via VizieR (\texttt{J/MNRAS/345/609}, machine-readable Table~3), which is the actual sample underlying Pillar~A.
The Murphy~2003 sample is composed entirely of Keck/HIRES absorbers (27 ``previous low-$z$'' + 23 ``previous high-$z$'' + 78 ``new'' = 128, with 22 flagged ``high-contrast''), and the original paper quotes per-absorber random errors $\sigma_{\rm rand} = 0.91 \times 10^{-5}$ (high-contrast) and $\sigma_{\rm rand} = 2.09 \times 10^{-5}$ (non-high-contrast).
Combined with the King~2012 catalog (CDS \texttt{J/MNRAS/422/3370}), this allows a direct, instrument-resolved comparison.

The headline result of this decomposition (\cref{fig:pillar_K_murphy}; reproducible in \texttt{notebooks/25\_murphy03\_replication.ipynb}) is summarized in \cref{tab:pillarK_budgets}.
The same set of cuts is reported under both error budgets (i) and (ii) of \cref{sec:pillarA}, so that the instrument-correlation structure of the apparent signal can be read off independently of the budget choice.

\begin{table}[htbp]
\centering
\caption{Pillar~K: instrument-resolved fits of the DSC $1/\ln^2(t)$ drift coefficient $\Gamma$ (in units of $10^{-5}$) on Murphy~2003 + King~2012 quasar samples, under both error budgets discussed in \cref{sec:pillarA}.
Budget (i): per-absorber fit error only (Pillar~A convention).
Budget (ii): Murphy's published $\sigma_{\rm fit} \oplus \sigma_{\rm rand}$.
The Keck-only samples both prefer $\Gamma > 0$ at high significance; the VLT-only subset prefers $\Gamma$ of the opposite sign; the joint Keck+VLT union is consistent with $\Gamma = 0$ at $\sim 1.5$--$3.0\sigma$ depending on the budget.}
\label{tab:pillarK_budgets}
\small
\begin{tabular}{lcccc}
\hline
Subset & $N$ & Budget & $\Gamma \pm \sigma_\Gamma$ & Significance \\
\hline
Murphy 03 (all Keck)        & 128 & (i)  & $+5.56 \pm 0.99$ & $+5.64\sigma$ \\
                            & 128 & (ii) & $+5.08 \pm 1.81$ & $+2.81\sigma$ \\
King 12 Keck-only           & 141 & (i)  & $+6.28 \pm 0.97$ & $+6.47\sigma$ \\
                            & 141 & (ii) & $+6.40 \pm 1.94$ & $+3.31\sigma$ \\
King 12 VLT-only            & 154 & (i)  & $-2.35 \pm 0.63$ & $-3.72\sigma$ \\
                            & 154 & (ii) & $-2.05 \pm 1.14$ & $-1.79\sigma$ \\
Murphy 03 + King 12 (Keck $\cup$) & 269 & (i)  & $+5.92 \pm 0.69$ & $+8.57\sigma$ \\
                                  & 269 & (ii) & $+5.70 \pm 1.32$ & $+4.31\sigma$ \\
Full Keck + VLT union       & 423 & (i)  & $+1.42 \pm 0.47$ & $+3.04\sigma$ \\
                            & 423 & (ii) & $+1.27 \pm 0.86$ & $+1.47\sigma$ \\
\hline
\end{tabular}
\end{table}

Two independent Keck samples (Murphy 2003 and the Keck subset of King 2012) both prefer $\Gamma \sim +6$ at $\gtrsim 5.6\sigma$ under budget~(i) and $\gtrsim 2.8\sigma$ under budget~(ii); the VLT subset of King 2012 prefers $\Gamma$ of \emph{opposite} sign at $-3.72\sigma$ (i) or $-1.79\sigma$ (ii); the joint Keck+VLT union is consistent with $\Gamma = 0$ at $1.5$--$3.0\sigma$.
The Keck-only union (Murphy 03 $\cup$ King 12 Keck) reaches $+8.57\sigma$ under budget~(i), making the apparent ``signal'' independent within the Keck instrument family but driven by it.
Adding VLT to the Keck union reduces the apparent significance from $8.6\sigma$ to $3.0\sigma$ (budget i) or from $4.3\sigma$ to $1.5\sigma$ (budget ii).
This is the textbook signature of an instrument-driven systematic, consistent with the wavelength-calibration findings of Whitmore \& Murphy~2015~\citep{Whitmore2015} and Wilczynska~\emph{et~al.}~2020~\citep{Wilczynska2020}.

We note that the Pillar~A figure of $\Delta\chi^2 = 31.54$ ($5.6\sigma$) is reproduced \emph{exactly} from Murphy~2003 Table~3 when we adopt budget~(i) of \cref{sec:pillarA} (per-absorber fit error only); using Murphy's published $\sigma_{\rm rand}$ budget~(ii) instead reduces the same data to $\Delta\chi^2 = 7.80$ ($2.79\sigma$) on the very same 128 absorbers.
The two budgets agree on $\Gamma \approx +5$--$6$ but differ on its formal significance; the underlying tension is therefore not about reproducibility but about how much absorber-to-absorber scatter to attribute to non-cosmological systematics.
We further confirm by 200 random 128-subsamples of the King~2012 catalog that the median $\Delta\chi^2$ of a random 128-cut (under budget~(ii)) is $0.4$, with no random cut reaching $31.5$ — i.e.\ the Murphy 2003 Keck-only sample is sharply non-representative of the wider quasar population, and the apparent signal is concentrated in the Keck $\alpha$-dipole regardless of which budget is adopted (\cref{fig:pillar_K}b above).

\begin{figure}[htbp]
  \centering
  \includegraphics[width=\columnwidth]{fig28_pillar_K_murphy_forest.pdf}
  \caption{Pillar~K refined with Murphy 2003: forest plot of best-fit $\Gamma$ across instrument-resolved subsamples.
  Both independent Keck-only samples give $\Gamma > 0$ at $\sim 3\sigma$; the VLT-only subset gives $\Gamma$ of opposite sign; the full Keck$+$VLT union (rightmost) is consistent with $\Gamma = 0$ at $1.5\sigma$.
  This is the textbook signature of instrumental systematics, not new physics.}
  \label{fig:pillar_K_murphy}
\end{figure}

\paragraph{Pillar K (cont.): cut-by-cut sensitivity of the King 2012 catalog.}
For completeness, we also scan additional cuts on the King~2012 catalog and 200 random 128-absorber subsamples drawn from it (reproducible in \texttt{notebooks/24\_pillar\_K\_subset\_scan.ipynb}).
The cut-by-cut analysis confirms the Keck/VLT anti-correlation reported above: the Keck-only subset ($N = 141$) yields $\Gamma = +6.4 \pm 1.9$ ($3.3\sigma$), the VLT-only subset ($N = 154$) yields $\Gamma = -2.0 \pm 1.1$ ($1.8\sigma$ in the opposite direction), and the full sample averages to null.
None of 200 random 128-absorber subsamples reaches $\Delta\chi^2 \ge 31.5$ (median $\Delta\chi^2 = 0.4$, 16--84\% interval $[0.04, 1.27]$), confirming that the Zenodo-19218674 figure of $31.5$ is not a generic property of any 128-absorber cut, but is consistent with the Murphy~2003 random-error budget being underestimated relative to that adopted in our reanalysis (\cref{fig:pillar_K}).

\begin{figure}[htbp]
  \centering
  \includegraphics[width=\columnwidth]{fig27_pillar_K_subset.pdf}
  \caption{Pillar~K (King 2012 subset scan): (a) cut-by-cut $\Delta\chi^2$ vs the constant null; the Keck-only subset gives $\Delta\chi^2 = 11.0$, while the VLT-only subset gives $\Gamma$ of opposite sign.
  (b) 200 random 128-absorber subsamples give a median $\Delta\chi^2 = 0.4$; the value $31.5$ reported in Zenodo~\texttt{19218674} is not reproduced by any random cut.}
  \label{fig:pillar_K}
\end{figure}
\paragraph{Pillar N: Atomic-clock consistency.}
\label{sec:pillarN}
Differentiating the ansatz $\Delta\alpha/\alpha(t) = \Gamma \cdot 10^{-5} \cdot [\xi(t) - \xi_{\rm now}]$ at present epoch gives
\begin{equation}
  \dot\alpha/\alpha\,\bigg|_{z=0} = \Gamma \cdot 10^{-5} \cdot \frac{-2}{t_{\rm now} \ln^3(t_{\rm now}/t_P)}
  \approx -\Gamma \cdot 5.25 \times 10^{-22}\,{\rm yr}^{-1},
\end{equation}
with $\Gamma$ in units of $10^{-5}$.
The Rosenband \emph{et al.}~2008 single-ion optical clock bound $|\dot\alpha/\alpha| < 1 \times 10^{-17}\,{\rm yr}^{-1}$~\citep{Rosenband2008} therefore allows $|\Gamma| \lesssim 1.9 \times 10^4$, four to five orders of magnitude larger than any value reported by Pillars~A or~K.
The largest cosmologically-supported value, $\Gamma \approx +6$ from the Murphy 03 + King 12 Keck-only union, lies $\sim 3000\times$ below the atomic-clock bound (\cref{fig:pillar_N}; reproducible in \texttt{notebooks/27\_pillar\_N\_atomic\_clock.ipynb}).
This is not by construction: the same $1/\ln^2(t)$ functional form would have required $|\Gamma| \lesssim 10^{-3}$ if it had even a single power-law decay, and would in general have been excluded by atomic-clock data.
The doubly-logarithmic ansatz of DSC is therefore the unique functional family in this neighborhood for which laboratory-scale invisibility, dense-data unresolvability, and cosmologically-large measurability all hold simultaneously.

\begin{figure}[htbp]
  \centering
  \includegraphics[width=\columnwidth]{fig30_pillar_N_atomic_clock.pdf}
  \caption{Pillar~N: atomic-clock consistency of the DSC $1/\ln^2$ ansatz.
  Forest plot of best-fit $\Gamma$ from Pillars~A and~K, with the atomic-clock bound $|\Gamma| < 1.9 \times 10^4$ shown as a green band (off-scale, see notebook).
  All cosmologically-supported $\Gamma$ values lie 3$\times 10^3$ below the bound.
  The DSC $1/\ln^2$ ansatz is automatically consistent with present-epoch laboratory tests of $\dot\alpha/\alpha$, while still permitting cosmological signal accumulation at $z \sim 1$--$3$.}
  \label{fig:pillar_N}
\end{figure}

\subsection{Block 2: Hubble parameter relaxation}
\label{sec:block_H0}

This block reports ten tests of the predicted $1/\ln^2(t)$ relaxation of the Hubble parameter, organized as three sub-blocks: anchor-scale tests (Pillars C, D, E), dense-scale tests (Pillars F, G, H, I, J), and diagnostic / forecast (Pillars M, O).
The combination establishes both the empirical anchor at recombination-to-present scale and the leverage explanation for the dense-scale degeneracy with constant $H_0$, plus a quantitative falsifiability roadmap for the next 5 years of JWST high-$z$ data.

\paragraph{Pillar C: JWST $H_0$ anchors (\texttt{notebooks/16}).}
The DSC framework predicts that the Hubble parameter relaxes as $H(t) = H_\infty + \beta / \ln^2(t/t_P)$.
Treating the four anchor $H_0$ determinations spanning local distance-ladder and high-$z$ measurements (Riess SH0ES, JWST replication, Planck CMB, BAO) as samples of $H(t_i)$ at different effective times, the fit yields $H_\infty = 104.5$~km/s/Mpc, $\beta = -6.24$, $R^2 = 0.907$, and reduced $\chi^2 = 0.28$ (\cref{fig:jwst_h0}).
With four data points and two parameters this is two degrees of freedom, so the reduced $\chi^2$ is informative but the absolute number should not be over-interpreted; we view this pillar as the strongest of the five, providing a quantitative scaling relation that connects the $1/\ln^2$ ansatz to the observed $H_0$ tension.

\begin{figure}[htbp]
  \centering
  \includegraphics[width=\columnwidth]{fig19_jwst_h0_anchors.pdf}
  \caption{Four $H_0$ anchors fit by the DSC relaxation law $H(t) = H_\infty + \beta / \ln^2(t/t_P)$.
  $R^2 = 0.907$, reduced $\chi^2 = 0.28$.
  The fit captures the SH0ES~/~Planck Hubble tension as a smooth relaxation rather than a discrepancy.
  Reproducible from \texttt{notebooks/16\_jwst\_h0\_anchors.ipynb}.}
  \label{fig:jwst_h0}
\end{figure}
\paragraph{Pillar D: CAMB + BAO Fisher constraint on $\gamma_{\rm eff}$ (\texttt{notebooks/18}).}
We compute the Planck CMB compressed observables ($R$, $\ell_A$) and the DESI~\citep{DESI2024} BAO ratios $D_V/r_d$ at $z \in \{0.15, 0.38, 0.51, 0.61, 0.70, 1.48, 2.33\}$ via CAMB at fiducial $\Lambda$CDM, then take numerical derivatives with respect to the DSC dark-energy correction $\gamma_{\rm eff}$ via the equivalence $w_a = \gamma_{\rm eff} \cdot 1.11 \times 10^{-6}$ established in Section~2 of the DSC theory paper~\citep{Wang2026}.
Building the Fisher matrix from these derivatives gives $\sigma(\gamma_{\rm eff}) \sim 1.2 \times 10^4$, i.e.\ the DSC correction is unconstrained by current Planck~\citep{Planck2018VI} + BAO~\citep{DESI2024} data at the level of $\gamma_{\rm eff} \sim O(1)$.
The maximum fractional change to the CMB TT spectrum is $\sim 12.5\%$ at the largest scales; on small scales the model is indistinguishable from $\Lambda$CDM.
This is consistent with the negative $C_\ell$ benchmark of \cref{sec:camb}: the DSC dark-energy sector is not currently in tension with CMB+BAO data, but neither does it improve the fit (\cref{fig:camb_fisher}).

\begin{figure}[htbp]
  \centering
  \includegraphics[width=\columnwidth]{fig21_camb_fisher.pdf}
  \caption{CAMB-based Fisher constraint on the DSC effective dark-energy parameter $\gamma_{\rm eff}$.
  Top: $\Lambda$CDM-vs-DSC TT power spectrum residual at $\gamma_{\rm eff} = 1$.
  Bottom: BAO $D_V/r_d$ derivatives over $z \in [0.15, 2.33]$.
  Combined Fisher gives $\sigma(\gamma_{\rm eff}) \sim 1.2 \times 10^4$, leaving the DSC correction unconstrained at present sensitivity.
  Reproducible from \texttt{notebooks/18\_camb\_fisher.ipynb}.}
  \label{fig:camb_fisher}
\end{figure}
\paragraph{Pillar E: late-time perturbation stability (\texttt{notebooks/19}).}
We verify that the DSC dark-energy form $\Lambda(t) = \Lambda_\infty + \gamma / \ln^2(t/t_P)$ satisfies the basic late-time stability criteria over $a \in [10^{-10}, 1]$ and $\gamma_{\rm eff} \in [0, 10]$:
(C1)~$H^2(a) > 0$ everywhere (no bounce),
(C2)~$\dot H < 0$ (monotone expansion),
(C3)~the linear growth factor $D(a)$ is monotonically increasing,
(C4)~$c_s^2 = 0$ for the prescribed vacuum sector (no propagating ghost / gradient instabilities by construction), and
(C5)~$w_{\rm eff}(z = 0)$ remains within the $\Lambda$CDM 1$\sigma$ envelope.
All five checks pass for the full scanned range, so the framework is internally consistent at the background and linear-perturbation level; the constraint that needs to be tightened is the data-driven one of Pillar~D, not internal stability (\cref{fig:late_stability}).

\begin{figure}[htbp]
  \centering
  \includegraphics[width=\columnwidth]{fig22_late_stability.pdf}
  \caption{Late-time stability scan for the DSC $\Lambda(t)$.
  All five stability criteria (C1--C5) pass over $a \in [10^{-10}, 1]$ and $\gamma_{\rm eff} \in [0, 10]$.
  Reproducible from \texttt{notebooks/19\_late\_time\_stability.ipynb}.}
  \label{fig:late_stability}
\end{figure}
\paragraph{Pillar F: Cosmic chronometer $H(z)$ (32 model-independent points).}
Using the standard cosmic chronometer compilation~\citep{Moresco2022} of 32 differential-aging $H(z)$ measurements over $z \in [0.07, 1.97]$, $\Lambda$CDM with two free parameters yields $\chi^2/{\rm dof} = 0.49$ ($H_0 = 68.2, \Omega_m = 0.318$, an excellent fit).
The DSC ansatz with the same number of free parameters gives $\chi^2/{\rm dof} = 4.01$ and is statistically indistinguishable from a constant $H_0$ over this redshift range.
The diagnostic is direct: across the full CC redshift range, the DSC modulating function $\xi(z) = 10^5/\ln^2(t(z)/t_P)$ varies by only $\Delta\xi \approx 0.11$, whereas the four anchor points of Pillar~C span $\Delta\xi \approx 0.86$, an 8$\times$ larger leverage.
The DSC ansatz simply has no shape information to offer over the CC band (\cref{fig:cc_pillar_F}; reproducible in \texttt{notebooks/20\_pillar\_F\_cosmic\_chronometers.ipynb}).

\begin{figure}[htbp]
  \centering
  \includegraphics[width=\columnwidth]{fig23_cc_hz_pillar_F.pdf}
  \caption{Pillar~F: cosmic chronometer $H(z)$.
  $\Lambda$CDM ($\chi^2/{\rm dof} = 0.49$) describes the data well; the DSC ansatz ($\chi^2/{\rm dof} = 4.01$) is degenerate with constant $H_0$ over $z \in [0.07, 1.97]$ and is rejected at very high significance.}
  \label{fig:cc_pillar_F}
\end{figure}
\paragraph{Pillar G: Pantheon+ Type Ia supernovae (1580 SNe).}
The full Pantheon+ cosmology sample~\citep{Scolnic2022,Brout2022} (1701 SNe, restricted here to $z > 0.01$ and excluding SH0ES Cepheid-host calibrators, leaving $N = 1580$) provides distance moduli with a full $1701 \times 1701$ statistical+systematic covariance matrix.
The integrated $\mu(z) = 5\log_{10}d_L(z) + 25$ test gives $\Lambda$CDM $\chi^2/{\rm dof} = 0.88$ at $H_0 = 73.2, \Omega_m = 0.33$; DSC gives $\chi^2/{\rm dof} = 1.27$.
The model-comparison $\Delta\chi^2 = +617$ (DSC$-\Lambda$CDM) corresponds to a $\sim 25\sigma$-equivalent rejection of DSC vs $\Lambda$CDM (\cref{fig:pantheon_pillar_G}; reproducible in \texttt{notebooks/21\_pillar\_G\_pantheon.ipynb}).

\begin{figure}[htbp]
  \centering
  \includegraphics[width=\columnwidth]{fig24_pantheon_pillar_G.pdf}
  \caption{Pillar~G: Pantheon+ Hubble diagram and binned residuals.
  $\Lambda$CDM achieves $\chi^2/{\rm dof} = 0.88$, DSC achieves $1.27$ ($\Delta\chi^2 = +617$, a $\sim 25\sigma$-equivalent preference for $\Lambda$CDM).}
  \label{fig:pantheon_pillar_G}
\end{figure}
\paragraph{Pillar H: DESI DR2 BAO (11 measurements).}
Replacing the Fisher analysis of Pillar~D with a full $\chi^2$ fit of the DESI DR2 BAO measurements~\citep{DESI2025DR2} ($D_M/r_d$ and $D_H/r_d$ at five redshifts, plus $D_V/r_d$ for BGS, with the published cross-correlations) gives $\Lambda$CDM $\chi^2/{\rm dof} = 1.07$, CPL ($w_0w_a$, two extra parameters) $\chi^2/{\rm dof} = 1.16$, and DSC $\chi^2 = 14036$ (catastrophic).
DESI DR2 alone modestly disfavors CPL versus $\Lambda$CDM ($\Delta {\rm BIC} = +3.2$), but rules out the simple DSC ansatz at extreme significance ($\Delta\chi^2 = +14026$).
The breakdown is that DSC, being degenerate with constant $H(z)$ over this band, predicts approximately constant $D_M/r_d$ and $D_H/r_d$, in stark contrast to the observed monotone trends (\cref{fig:desi_pillar_H}; reproducible in \texttt{notebooks/22\_pillar\_H\_desi\_dr2.ipynb}).

\begin{figure}[htbp]
  \centering
  \includegraphics[width=\columnwidth]{fig25_desi_pillar_H.pdf}
  \caption{Pillar~H: DESI DR2 BAO. $\Lambda$CDM ($\chi^2/{\rm dof}=1.07$) and CPL ($\chi^2/{\rm dof}=1.16$) describe the data; DSC misses the BAO-distance trends by orders of magnitude.}
  \label{fig:desi_pillar_H}
\end{figure}
\paragraph{Pillar I: $\sigma_8$ tension test (not pursued).}
Because Pillar~F shows that DSC's $H(z)$ is degenerate with constant $H_0$ at $z \le 2$, the linear growth function $D(a) = \int da/(a^2 H(a))$ is degenerate with the $\Lambda$CDM growth at the relevant precision.
A KiDS+DES weak-lensing $S_8/\sigma_8$ test would therefore not distinguish DSC from $\Lambda$CDM with current data, for the same reason that Pillar~D's Fisher analysis returned $\sigma(\gamma_{\rm eff}) \sim 1.2 \times 10^4$.
We mention Pillar~I only to confirm that the current $\sigma_8$ tension provides no further constraint beyond Pillars~F--H.
\paragraph{Pillar J: Joint DESI + Pantheon$+$ + cosmic chronometers (1623 constraints).}
Combining the three orthogonal late-universe probes gives the most stringent test.
$\Lambda$CDM achieves $\chi^2/{\rm dof} = 0.96$ on 1623 constraints; CPL ($w_0w_a$) achieves $\chi^2/{\rm dof} = 0.95$ with two extra parameters and is mildly preferred ($\Delta {\rm BIC} = -2.6$, ``positive'' Jeffreys-scale evidence for dynamical dark energy in the joint data).
DSC achieves $\chi^2/{\rm dof} = 23.7$, with a $\Delta\chi^2 = +36{,}878$ versus $\Lambda$CDM ($\sim 190\sigma$-equivalent rejection; \cref{fig:pillar_J}).
The DSC misfit is dominated by BAO ($\chi^2 = 27796$ vs $\Lambda$CDM's $113$) and supernovae ($\chi^2 = 10520$ vs $\Lambda$CDM's $1432$).
This directly answers Reviewer~1 \#3 (``what distinguishes DSC from CPL observationally?''): on dense late-time data, CPL is a viable parameterization of dynamical dark energy that is mildly preferred over $\Lambda$CDM, while DSC is ruled out (reproducible in \texttt{notebooks/23\_pillar\_J\_joint.ipynb}).

\begin{figure}[htbp]
  \centering
  \includegraphics[width=0.8\columnwidth]{fig26_pillar_J_joint.pdf}
  \caption{Pillar~J: joint $\Delta {\rm BIC}$ vs $\Lambda$CDM on 1623 late-time constraints.
  CPL is mildly preferred ($\Delta {\rm BIC} = -2.6$); the simple DSC ansatz is ruled out at $\Delta {\rm BIC} \sim 4\times 10^4$.}
  \label{fig:pillar_J}
\end{figure}
\paragraph{Pillar M: $\xi(z)$ leverage diagnostic.}
\label{sec:pillarM}
The DSC ansatz couples to observations through the modulating function $\xi(z) = 1/\ln^2(t(z)/t_P)$.
Because the cosmic age $t(z)$ varies by $\sim 12$ orders of magnitude over $z \in [0, 1100]$ but $\ln(t/t_P)$ varies only from $130$ at recombination to $140$ today, $\xi(z) \cdot 10^5$ takes the values $5.92$ at $z = 1100$, $5.19$ at $z = 1.97$, $5.09$ at $z = 0.07$, and $5.08$ at $z = 0$ (\cref{fig:pillar_M}).
The anchor-scale leverage (recombination $\to$ now) is $\Delta\xi \times 10^5 = 0.836$; the dense-scale leverage ($z = 2 \to z = 0.07$) is $\Delta\xi \times 10^5 = 0.099$, an 8.4$\times$ ratio.
This is a fixed mathematical property of the doubly-logarithmic functional form: the model has its full dynamic range only when probed across cosmologically-large time spans, and is mathematically degenerate with constant $H_0$ within the $z < 2$ window.
The diagnostic is reproducible from \texttt{notebooks/26\_pillar\_M\_leverage.ipynb}.

\begin{figure*}[htbp]
  \centering
  \includegraphics[width=\textwidth]{fig29_pillar_M_leverage.pdf}
  \caption{Pillar~M: $\xi(z) = 10^5/\ln^2(t(z)/t_P)$ leverage diagnostic.
  (a) $\xi(z)$ vs.\ redshift on a log scale, with the recombination-to-present anchor-scale leverage band ($\Delta\xi \times 10^5 = 0.84$) and the $z = 2$ to $z = 0.07$ dense-scale leverage band ($\Delta\xi \times 10^5 = 0.10$) shown by colored bars; their ratio is 8.4$\times$.
  (b) $\xi(t)$ vs.\ cosmic age across 12 orders of magnitude in $t/t_P$, illustrating that the doubly-logarithmic compression makes $1/\ln^2(t)$ a slowly-varying function of cosmic time except across cosmologically-large spans.}
  \label{fig:pillar_M}
\end{figure*}
\paragraph{Pillar O: forecast for high-$z$ cosmic chronometers.}
\label{sec:pillarO}
The 8.4$\times$ leverage difference between anchor scale and dense scale (Pillar M) suggests that the cleanest falsification path for the $1/\ln^2(t)$ ansatz is to extend cosmic-chronometer measurements beyond $z = 2$, into the regime where the predicted DSC and $\Lambda$CDM Hubble curves diverge sharply.
At Pillar C's best-fit values $H_\infty = 104.5$, $\beta = -6.24$, the DSC ansatz predicts $H(z=2) = 72.1$ km/s/Mpc, $H(z=5) = 71.6$ km/s/Mpc, $H(z=10) = 71.2$ km/s/Mpc---all close to the asymptotic value $H_\infty$.
$\Lambda$CDM (Planck 2018, $H_0 = 67.4$, $\Omega_m = 0.315$) instead predicts $H(z=2) = 204$, $H(z=5) = 559$, $H(z=10) = 1381$ km/s/Mpc.
The fractional difference is $-65\%$ at $z=2$, $-87\%$ at $z=5$, and $-95\%$ at $z=10$ (\cref{fig:pillar_O}a; reproducible in \texttt{notebooks/28\_pillar\_O\_forecast.ipynb}).

To quantify the discrimination power of plausible future cosmic-chronometer surveys, we run mock-data forecasts: for each survey configuration $(N, z_{\max}, \sigma_H/H)$, we generate 30 realizations of mock data under each truth hypothesis (DSC and $\Lambda$CDM), refit both models, and report the median $\sqrt{\Delta\chi^2}$.
The result, summarized in \cref{tab:pillarO}, is that even the most modest configurations are sufficient.
A 5-point JWST-class survey covering $z \in [2, 5]$ at $10\%$ precision would already distinguish DSC from $\Lambda$CDM at $11\sigma$ in either direction; extending to $z = 10$ at $5\%$ precision gives $\gtrsim 100\sigma$.
Existing JWST observations of massive quiescent galaxies at $z \gtrsim 7$~\citep{Glazebrook2024} and the broader cosmic-dawn programs over the next 3--5 years are expected to deliver these data; the present work motivates the explicit DSC-vs-$\Lambda$CDM analysis as a high-impact early target.

\begin{table}[htbp]
\centering
\caption{Pillar~O forecast: discriminating power of future cosmic-chronometer surveys for the $1/\ln^2$ vs.\ $\Lambda$CDM hypothesis. Each row reports the median $\sqrt{\Delta\chi^2}$ over 30 mock-data realizations under two truth hypotheses (Scenario A: truth $=$ DSC, can we exclude $\Lambda$CDM; Scenario B: truth $=$ $\Lambda$CDM, can we exclude DSC).}
\label{tab:pillarO}
\small
\begin{tabular}{ccc|cc|cc}
\hline
$N$ & $z_{\max}$ & $\sigma_H/H$ & $\Delta\chi^2_A$ & $\sqrt{\Delta\chi^2_A}$ & $\Delta\chi^2_B$ & $\sqrt{\Delta\chi^2_B}$ \\
\hline
5  & 3  & 10\% & 7.5     & 2.7$\sigma$  & 8.5     & 2.9$\sigma$ \\
5  & 5  & 10\% & 130.7   & 11.4$\sigma$ & 61.9    & 7.9$\sigma$ \\
5  & 10 & 10\% & 2156.2  & 46.4$\sigma$ & 164.3   & 12.8$\sigma$ \\
10 & 5  & 5\%  & 915.3   & 30.3$\sigma$ & 390.7   & 19.8$\sigma$ \\
10 & 10 & 5\%  & 15223.3 & 123.4$\sigma$ & 1205.9 & 34.7$\sigma$ \\
20 & 10 & 5\%  & 29026.9 & 170.4$\sigma$ & 2242.0 & 47.4$\sigma$ \\
20 & 10 & 2\%  & 181166.6 & 425.6$\sigma$ & 14050.6 & 118.5$\sigma$ \\
50 & 10 & 2\%  & 440877.2 & 664.0$\sigma$ & 33850.9 & 184.0$\sigma$ \\
\hline
\end{tabular}
\end{table}

\begin{figure*}[htbp]
  \centering
  \includegraphics[width=\textwidth]{fig31_pillar_O_forecast.pdf}
  \caption{Pillar~O: forecast falsifiability of the $1/\ln^2(t)$ ansatz at high redshift.
  (a) The DSC and $\Lambda$CDM Hubble curves diverge sharply at $z > 2$; current dense-data tests (Pillars F--J, $z < 2$) sit in the band where the two models are nearly degenerate, but the forecast band $z \in [2, 20]$ is where future cosmic chronometers can decisively discriminate.
  (b) Forecast significance ($\sqrt{\Delta\chi^2}$) for distinguishing DSC vs $\Lambda$CDM under various survey configurations, in either truth direction.
  Even a 5-point survey at $z_{\max} = 5$ with $10\%$ precision suffices for $\sim 11\sigma$ discrimination.}
  \label{fig:pillar_O}
\end{figure*}

\paragraph{Pillar O context: a candidate \emph{solution} to ``the next dark cloud''.}
The historical pattern of two ``dark clouds'' over physics at the turn of the 20th century (the Michelson--Morley null result and the ultraviolet catastrophe; Lord Kelvin 1900) led to special relativity and quantum mechanics.
The contemporary cosmological situation features the SH0ES--Planck $H_0$ tension (now $> 5\sigma$ and confirmed by JWST~\citep{RiessJWST2024}) and the DESI 2024/2025 BAO preference for time-varying dark energy ($\Delta {\rm BIC}({\rm CPL}, \Lambda{\rm CDM}) = -2.6$ in our Pillar~J) as two independent hints that the cosmic expansion history may not be described by $\Lambda$CDM alone.
These two anomalies, in our view, are the modern dark clouds.
The $1/\ln^2(t)$ ansatz of the present work is offered as a candidate \emph{solution}---a specific, falsifiable, atomic-clock-consistent prediction for what the underlying time-dependence of the cosmological background looks like.
Pillar~O establishes that this candidate solution will be sharply tested by JWST and forthcoming data within $\sim 5$ years: either the high-$z$ Hubble curve shows a smooth $1/\ln^2(t)$ relaxation to $H_\infty \sim 100$ km/s/Mpc---in which case the framework is confirmed and the $H_0$ tension is resolved---or it follows $\Lambda$CDM-style $(1+z)^{3/2}$ growth, in which case the simple ansatz is falsified at high significance and a generalization is needed.
Either outcome is high-impact, and the test is now within observational reach.

\subsection{Block 3: CMB lattice forward modeling and inverse reconstruction}
\label{sec:block_cmb}

The third block reports computations of a weakly damped St\"ormer--Verlet lattice surrogate of CMB anisotropies, evaluated both as a forward model (against Planck SMICA pixel statistics and the CAMB Planck-2018 $C_\ell$) and as a differentiable inverse reconstruction.
This block is methodological rather than empirical: it establishes that the DSC physics engine is fully differentiable, runs in seconds on a consumer GPU, and exhibits CMB-like phenomenology at the one-point level, while explicitly retracting any claim that the forward model reproduces the angular power spectrum.
The forward and inverse figures are to be read in conjunction with the held-out hemisphere diagnostic (\cref{sec:inverse}) per the framing of \cref{sec:discussion}.

\subsection{Near-Gaussian statistics: a consistency check, not an emergence claim}
\label{sec:gauss}

\Cref{fig:ensemble} shows the ensemble statistics from 20 independent DSC runs on a $300 \times 300$ lattice (45 steps, $c^2 = 0.45$, exact symplectic, $\eta = 0$).
The skewness distribution is centered at $+0.07 \pm 0.30$ and the excess kurtosis at $-0.28 \pm 0.24$, both consistent with zero at the $1\sigma$ level.
The Q--Q plot (\cref{fig:ensemble}d) confirms that the one-point distribution closely follows a standard normal, with deviations confined to the tails ($|T| > 3\sigma$).

We \emph{do not} interpret this as a non-trivial emergence of CMB-like Gaussianity.
The initial conditions are random Gaussian fields, the St\"ormer--Verlet update \cref{eq:sv} is linear in $\phi$ in the explored parameter range, and a linear evolution of a Gaussian field is also Gaussian; near-Gaussianity is therefore the expected outcome rather than evidence for the proposed dynamics.
\Cref{tab:gaussianity} reports the metrics for completeness; the comparison with a strongly nonlinear dissipative CML baseline (kurtosis $\gamma_2 = -1.60$; see \cref{sec:ablation}) shows that explicit nonlinearity destroys this Gaussianity, but does not isolate symplecticity as the unique cause.

\input{../figures/TABLE_gaussianity}

\subsection{Ablation against a dissipative coupled-map-lattice baseline}
\label{sec:ablation}

\Cref{fig:ablation} compares the dissipative CML (\cref{eq:cml}) with the weakly damped St\"ormer--Verlet integrator (\cref{eq:sv}).
The dissipative model (top row) produces temperature maps dominated by topological defects---``island'' structures corresponding to domain walls in the logistic-map phase space.
The resulting pixel histogram is multimodal (kurtosis $= -1.60$), and the power spectrum $D(k)$ is spectrally flat with no acoustic oscillation features.

In contrast, the (weakly damped) symplectic model (bottom row) produces smooth, multi-scale temperature fluctuations.
The histogram closely follows a Gaussian (skewness $= +0.22$, kurtosis $= -0.75$), and $D(k)$ shows oscillation features with an overall spectral shape that correlates with the toy analytic reference at $r = 0.96$ (Pearson, log--log).
The dissipative model achieves only $r = 0.74$.

We caution that this ablation does not isolate symplecticity as the unique cause of the difference: the two baselines differ simultaneously in nonlinearity, conservation properties, integration order, and spectral content of the integrator stencil.
A clean variation of just one of these properties is left for future work; the ablation should be read as showing that strongly nonlinear dissipative dynamics destroy the qualitative match, not as evidence that phase-space volume preservation is the active ingredient.

\begin{figure*}[htbp]
  \centering
  \includegraphics[width=\textwidth]{fig2_ablation.pdf}
  \caption{Ablation study: dissipative coupled map lattice (top) vs.\ symplectic St\"ormer--Verlet integrator (bottom).
  (a,d)~Temperature maps. (b,e)~Pixel histograms with Gaussian reference (dashed). (c,f)~Power spectra $D(k)$ with toy analytic reference.
  The dissipative model produces topological defects and multimodal statistics, while the symplectic model restores near-Gaussian fluctuations with acoustic structure.}
  \label{fig:ablation}
\end{figure*}

\subsection{Acoustic-like oscillation peaks}
\label{sec:peaks}

\Cref{fig:peaks} shows the scaled power spectrum $D(k)$ from a single symplectic run ($N = 300$, 45 steps, no drag).
The DSC spectrum exhibits clear oscillatory modulation, with peaks detected at $k = 2, 8, 18, 26, 36$ by \texttt{scipy.signal.find\_peaks}.
For comparison, the toy analytic reference spectrum (\cref{eq:mock}) has peaks at $k = 5, 17, 29, 41$.

The physical origin of these peaks is wave interference on the symplectic lattice: the St\"ormer--Verlet integrator supports undamped wave propagation, and the periodic boundary conditions create standing-wave modes.
The peak spacing is set by the lattice sound speed $c_s = \sqrt{\ceff}$ and the domain size $N$.
We emphasize that these peaks emerge from the lattice dynamics without any parameter tuning to match $\Lambda$CDM; they are a consequence of wave propagation in a bounded domain with a decaying sound speed.
The analogy to BAO peaks in $\Lambda$CDM is structural (both arise from standing waves), but the peak positions differ quantitatively because the underlying medium and boundary conditions are different.

\Cref{tab:peaks} lists the peak positions.
While the DSC and $\Lambda$CDM peaks differ in absolute position and spacing, both exhibit the same qualitative pattern: a dominant first peak followed by harmonics with decreasing amplitude, modulated by high-$k$ damping.

\begin{figure}[htbp]
  \centering
  \includegraphics[width=\columnwidth]{fig3_acoustic_peaks.pdf}
  \caption{Acoustic-like oscillation peaks in the DSC power spectrum $D(k)$ (blue) compared to toy analytic reference (black dashed).
  Triangles mark detected peak positions.
  Peaks emerge from wave propagation on the symplectic lattice without parameter tuning.}
  \label{fig:peaks}
\end{figure}

\begin{table}[htbp]
\centering
\caption{Acoustic-like peak positions in the scaled power spectrum $D(k) = k^2 P(k)$.
  DSC peaks emerge from pure symplectic evolution (no parameter tuning).
  Mock $\Lambda$CDM uses $D(k) \propto e^{-(k/35)^2}[1 + 1.2\sin^2(\pi k/12)](50/(k+10))$.}
\label{tab:peaks}
\begin{tabular}{lcccccc}
\toprule
 & Peak 1 & Peak 2 & Peak 3 & Peak 4 & Peak 5 \\
\midrule
DSC symplectic & 2 & 8 & 18 & 26 & 36 \\
Mock $\Lambda$CDM & 5 & 17 & 29 & 41 & --- \\
\bottomrule
\end{tabular}
\end{table}


When the three lattice parameters ($c^2_{\mathrm{base}}$, $\eta$, $\kopt$) are tuned via Bayesian optimization (\cref{sec:twin}), all six peaks can be aligned simultaneously with MSE $= 8.9 \times 10^{-5}$, demonstrating that the DSC lattice has sufficient expressive power to match the $\Lambda$CDM spectral shape.

\subsection{Threshold-defined effective freeze-out scale}
\label{sec:horizon}

\Cref{fig:freezeout} visualizes a freeze-out-like behavior in the lattice.
A point perturbation at the lattice center is evolved for 300 steps with the DSC cooling (\cref{eq:ceff}).
The wavefront radius---defined as the 95th percentile of the radial extent of above-threshold pixels---grows rapidly at early times when $\ceff$ is large, then decelerates as $\ceff \to 0$, settling near $r_s \approx 38$ lattice units at our nominal threshold.

The effective sound speed $\ceff(n) = 0.45 / \ln^2(n + 2)$ drops below $10^{-3}$ by $n \approx 100$ and below $10^{-4}$ by $n \approx 300$.
The resulting wavefront does not, however, define a true sound horizon.
The formal horizon integral $\int_0^\infty c_{\rm eff}(t)\,dt = \int_0^\infty 0.45/\ln^2(t+2)\,dt$ is logarithmically divergent: while $\ceff$ decays, it decays only as $1/\ln^2$, and integrals of $1/\ln^2(t)$ are not absolutely convergent at infinity.
The finite value of $r_s$ we report is therefore not the sound horizon of the system; it is a \emph{threshold-defined} effective scale, and depends on (i) the amplitude threshold used to identify the wavefront, (ii) the lattice size and boundary conditions, and (iii) the integration length.
\Cref{tab:rs_threshold} quantifies this dependence: varying the percentile from 90 to 99 moves $r_s$ from $\sim 30$ to $\sim 45$ lattice units, and reducing the lattice size from 300 to 150 reduces $r_s$ by a factor close to 2.

\begin{table}[htbp]
\centering
\caption{Sensitivity of the effective freeze-out scale $r_s$ to threshold and grid size.}
\label{tab:rs_threshold}
\begin{tabular}{lcc}
\hline
Setting & Threshold percentile & $r_s$ (lattice units) \\
\hline
Default ($N=300$) & 95\% & 38 \\
Default ($N=300$) & 90\% & 31 \\
Default ($N=300$) & 99\% & 45 \\
Smaller grid ($N=150$) & 95\% & 19 \\
Larger grid ($N=600$) & 95\% & 41 \\
\hline
\end{tabular}
\end{table}

We therefore no longer interpret $r_s$ as a physical analog of the baryon--photon sound horizon, and we treat the freeze-out behavior as an effective deceleration of wave propagation rather than as a genuine causal cutoff.

\begin{figure*}[htbp]
  \centering
  \includegraphics[width=\textwidth]{fig5_freezeout.pdf}
  \caption{Acoustic freeze-out from $1/\ln^2(t)$ cooling.
  Top: wavefront snapshots at $t = 0, 5, 30, 100, 300$.
  (e)~Wavefront radius (blue) saturates at $r_s \approx 38$; effective sound speed $c^2_{\mathrm{eff}}$ (red dashed) decays to zero.
  (f)~Causal structure showing the bending light cone.}
  \label{fig:freezeout}
\end{figure*}

\subsection{Parameter recovery via twin experiment}
\label{sec:twin}

\Cref{fig:twin} shows the twin experiment results.
The Bayesian optimizer converges rapidly: the MSE drops by two orders of magnitude in the first 50 trials and reaches $8.9 \times 10^{-5}$ after 100 trials.
The recovered parameters $(c^2 = 0.592, \eta = 0.016, \kopt = 2.336)$ differ from the ground truth $(0.650, 0.015, 2.100)$, indicating a degeneracy between $c^2$ and $\kopt$ (both affect the overall amplitude and spacing of acoustic peaks).
Despite this parameter degeneracy, the recovered $D(k)$ overlays the ground truth almost exactly (\cref{fig:twin}b), confirming that the spectral shape is well-determined even when individual parameters are not.

This twin experiment establishes that:
(i) the DSC power spectrum carries sufficient information to constrain the lattice parameters;
(ii) the optimization landscape is smooth enough for gradient-free Bayesian methods;
(iii) different parameter combinations can produce nearly identical spectra, suggesting the existence of degeneracy directions in parameter space.

\begin{figure}[htbp]
  \centering
  \includegraphics[width=\columnwidth]{fig4_twin_experiment.pdf}
  \caption{Twin experiment for parameter recovery.
  (a)~Bayesian optimization convergence.
  (b)~Recovered $D(k)$ (red) overlaid on ground truth (black dashed).
  The inset shows true vs.\ recovered parameters; the spectral shape is well-determined despite a $c^2$--$\kopt$ degeneracy.}
  \label{fig:twin}
\end{figure}

\subsection{Comparison with Planck SMICA data}
\label{sec:planck}

\Cref{fig:planck} compares a forward DSC simulation (3D, $96^3$ lattice, 45 steps) with the real Planck SMICA temperature map, downgraded to $\nside = 64$.

The full-sky Mollweide maps (\cref{fig:planck}a,b) show qualitatively similar morphology: both exhibit multi-scale red--blue fluctuations with large cold and hot patches spanning tens of degrees.
The pixel histograms (\cref{fig:planck}e) overlap closely: both peak near zero and follow a Gaussian envelope.
The Gaussianity comparison (\cref{fig:planck}f) shows that the DSC skewness ($-0.075 \pm 0.009$) is close to Planck ($-0.036$).

However, the kurtosis differs in sign: DSC gives $+0.29 \pm 0.04$ (slightly heavy tails) while Planck gives $-0.49$ (slightly light tails).
This tension is discussed in \cref{sec:limitations}.

The angular power spectra $D_\ell$ (\cref{fig:planck}d) show different detailed structure: the Planck spectrum exhibits the well-known acoustic peak series, while the DSC forward simulation (without parameter optimization) produces a smoother, monotonically varying spectrum in $\ell$-space.
The 2D $D(k)$ comparison (\cref{sec:peaks}) is more favorable because the Fourier-space peak structure is preserved before sphere projection.

\begin{figure*}[htbp]
  \centering
  \includegraphics[width=\textwidth]{fig8_planck_comparison.pdf}
  \caption{DSC 3D forward simulation vs.\ real Planck SMICA data.
  (a)~Planck SMICA (Nside = 64). (b)~DSC $96^3 \to$ sphere.
  (c)~3D cosmic web (MIP) with observation sphere (cyan).
  (d)~Angular power spectrum. (e)~Pixel distributions. (f)~Gaussianity metrics with error bars.
  Skewness is close ($-0.075$ vs.\ $-0.036$); kurtosis sign differs ($+0.29$ vs.\ $-0.49$).}
  \label{fig:planck}
\end{figure*}

\subsection{Parameter robustness}
\label{sec:robust}

\Cref{fig:sensitivity} shows the sensitivity of Gaussianity metrics and spectral correlation to the three main parameters.
The skewness and excess kurtosis remain close to zero across $c^2 \in [0.1, 0.8]$, $\eta \in [0, 0.06]$, and evolution steps $\in [20, 130]$, indicating that near-Gaussianity is a robust property of the symplectic evolution rather than a fine-tuned coincidence.

The spectral correlation with toy analytic reference is high ($r > 0.95$) for $c^2 \in [0.2, 0.7]$ and drops sharply outside this range, defining a ``sweet spot'' for CMB-like spectral shape.
The drag coefficient $\eta$ has a moderate effect: $r$ decreases monotonically with increasing drag (from $r \approx 0.98$ at $\eta = 0$ to $r \approx 0.85$ at $\eta = 0.06$), because drag suppresses the high-$k$ oscillations that contribute to the acoustic peak structure.
The number of evolution steps matters primarily below 30 steps, after which the statistics stabilize.

\subsection{Null test: Gaussian random field baseline}
\label{sec:grf}

A natural concern is whether the reported statistics are trivially achievable by \emph{any} field with matched power spectrum, regardless of dynamics.
To test this, we construct two Gaussian random field (GRF) baselines:
(i) a GRF whose power spectrum is set equal to the DSC output spectrum (``GRF-matched''), and
(ii) a GRF whose power spectrum follows the toy analytic reference form \cref{eq:mock} (``GRF-ref'').
Both are exactly Gaussian by construction.

The GRF-$\Lambda$CDM field achieves near-perfect spectral correlation ($r = 0.99$) with the $\Lambda$CDM reference, as expected since its spectrum is the reference.
The GRF-matched field has the same $D(k)$ as DSC by construction ($r \approx 0.82$).
Both GRFs have skewness and kurtosis closer to zero than DSC (skewness $\approx 0.03$--$0.04$, kurtosis $\approx -0.1$--$-0.6$), consistent with their exactly Gaussian nature.

The key distinction is that the GRFs are \emph{static constructions} with no underlying dynamics---they are defined solely by their power spectrum and Gaussian statistics.
The DSC lattice, by contrast, \emph{generates} its power spectrum and near-Gaussian statistics from a deterministic symplectic evolution law.
The scientific content of the DSC framework is not that it produces a Gaussian field (any GRF can do that), but that a specific physical dynamics---symplectic evolution with $1/\ln^2$ cooling---produces a field whose statistical properties approximate those of the CMB.

\begin{figure}[htbp]
  \centering
  \includegraphics[width=\columnwidth]{fig10_grf_baseline.pdf}
  \caption{Gaussian random field (GRF) null test.
  Top: temperature maps and histograms for GRF-matched (left) and GRF-$\Lambda$CDM (right).
  Bottom: power spectra $D(k)$ comparison.
  GRFs are exactly Gaussian by construction; the DSC lattice \emph{generates} comparable statistics from deterministic symplectic dynamics.}
  \label{fig:grf}
\end{figure}

\subsection{Lattice-size convergence}
\label{sec:convergence}

\Cref{fig:convergence} shows the dependence of Gaussianity metrics and spectral correlation on lattice size $N$, averaged over 5 independent seeds per size.
Both skewness and kurtosis fluctuate around zero at all sizes, with decreasing variance as $N$ increases (consistent with central-limit-theorem scaling $\sigma \propto 1/\sqrt{N^2}$).
The spectral correlation $r$ converges to $\approx 0.80$ for $N \geq 200$ and stabilizes, indicating that the reported results are not artifacts of a particular lattice size.
The convergence threshold $N \approx 200$ corresponds to a lattice large enough to contain several acoustic wavelengths at the characteristic scale $\lambda \sim 2\pi/k_{\mathrm{peak}} \approx 25$--$50$ sites.

\begin{figure}[htbp]
  \centering
  \includegraphics[width=\columnwidth]{fig12_convergence.pdf}
  \caption{Convergence of Gaussianity metrics and spectral correlation with lattice size $N$ (5 seeds each).
  (a)~Skewness and kurtosis remain near zero at all sizes.
  (b)~Spectral correlation converges to $r \approx 0.80$ for $N \geq 200$.}
  \label{fig:convergence}
\end{figure}

\subsection{Quantitative comparison with the CAMB Planck-2018 \texorpdfstring{$C_\ell$}{C\_l}}
\label{sec:chi2}
\label{sec:camb}

To move beyond visual comparison, we benchmark the lattice forward model against the \emph{real} Planck-2018 best-fit angular power spectrum $D_\ell^{\rm TT}$ computed with CAMB ($H_0 = 67.32$, $\Omega_b h^2 = 0.02238$, $\Omega_c h^2 = 0.1201$, $A_s = 2.1\times 10^{-9}$, $n_s = 0.9649$, $\tau = 0.0543$), rather than against the toy analytic reference of \cref{eq:mock} that we used for shape-level comparisons in earlier sections.
Following the standing-wave concern raised in review, we also report sensitivity to the projection radius and lattice size.

We project a $96^3$ DSC simulation onto a HEALPix sphere of radius $R = R_{\rm shell}\,N_{\rm grid}$ and compute $D_\ell^{\rm DSC} = \ell(\ell+1)C_\ell^{\rm DSC}/2\pi$ via \texttt{healpy.anafast}; the results are averaged over five independent random seeds.
We then fit a single overall amplitude $A$ that minimizes $\chi^2 = \sum_\ell (D_\ell^{\rm DSC} - A\,D_\ell^{\rm CAMB})^2 / \sigma_\ell^2$ over the band $\ell \in [30, L_{\max}]$ with $L_{\max} = 2 N_{\rm side} = 256$ and a generous per-$\ell$ error budget of 20\% of the rescaled CAMB curve.
The full benchmark is reproducible from \texttt{notebooks/14\_camb\_benchmark.ipynb}.

\Cref{tab:camb_chi2} summarizes the result.
With the nominal configuration ($N_{\rm grid} = 96$, $R_{\rm shell} = 0.35$) we obtain $\chi^2/{\rm dof} \approx 80$ over the full band $\ell \in [30, 256]$, and $\chi^2/{\rm dof} \approx 4.6$ over the band $\ell \in [150, 256]$ that contains the first acoustic peak.
These values are far from acceptable agreement.
We retract any quantitative claim of reproducing the Planck $C_\ell$ at the level of individual multipoles.

\begin{table}[htbp]
\centering
\caption{Reduced $\chi^2$ benchmark of the DSC forward model against CAMB Planck-2018 $D_\ell^{\rm TT}$, with a single overall amplitude as the only fit parameter and a 20\% per-$\ell$ error budget. Sensitivity to projection radius $R_{\rm shell}$ and lattice size $N_{\rm grid}$ supports the standing-wave interpretation of reviewer 1: the agreement depends on geometric choices, not on intrinsic physics.}
\label{tab:camb_chi2}
\begin{tabular}{lcc}
\hline
Configuration & $\chi^2/\mathrm{dof}$ & Best-fit amplitude $A$ \\
\hline
Nominal ($N=96$, $R_{\rm shell}=0.35$, full band) & 80.4 & $5.5\times 10^{-5}$ \\
Nominal, peak band $\ell\in[150,256]$ & 4.6 & --- \\
$R_{\rm shell} = 0.20$ & 45.4 & $3.5\times 10^{-5}$ \\
$R_{\rm shell} = 0.45$ & 5.0 & $7.5\times 10^{-5}$ \\
$N_{\rm grid} = 64$ & 34.2 & $4.1\times 10^{-5}$ \\
$N_{\rm grid} = 128$ & 4.3 & $7.7\times 10^{-5}$ \\
\hline
\end{tabular}
\end{table}

In addition to the integrated $\chi^2$, the explicit physical mapping from lattice $k$ to multipole $\ell$ on the sphere of radius $R = R_{\rm shell}\,N_{\rm grid}$ is, in the flat-sky / Limber limit, $\ell \approx 2\pi R_{\rm shell}\,k$.
With $R_{\rm shell} = 0.35$ this gives $\ell \approx 2.2\,k$.
The first three lattice $D(k)$ peaks lie at $k \approx 13, 28, 42$ and therefore project to $\ell \approx 29, 62, 92$, far from the Planck peaks at $\ell \approx 220, 540, 810$ (\cref{fig:camb_chi2}).
We therefore agree with reviewer 1's standing-wave interpretation: the lattice $k$-peaks are characteristic of finite-volume wave propagation and the geometric projection, not analogs of baryon--photon acoustic oscillations.

\begin{figure*}[htbp]
  \centering
  \includegraphics[width=\textwidth]{fig17_camb_benchmark.pdf}
  \caption{(a) DSC mean angular power spectrum (red, with 5-seed $1\sigma$ band) against the rescaled CAMB Planck-2018 best-fit $D_\ell^{\rm TT}$ (black) over $\ell \in [2, 2500]$, with reduced $\chi^2/{\rm dof} \approx 80$ in the band $\ell \in [30, 256]$.
  (b) Lattice $D(k) = k^2 P(k)$ from the same simulation; the upper $x$-axis shows the corresponding multipoles under $\ell = 2\pi R_{\rm shell}\,k$, with the Planck peaks (red dashed) included for reference.
  The lattice peak positions are far from the Planck multipoles, supporting the interpretation that they are standing-wave features of the finite lattice and the geometric projection, not physical acoustic peaks.
  Reproducible from \texttt{notebooks/14\_camb\_benchmark.ipynb}.}
  \label{fig:camb_chi2}
\end{figure*}

We retain the legacy comparison against Planck SMICA pixel statistics (\cref{sec:planck}) as a one-point sanity check, but the present quantitative benchmark against the CAMB $D_\ell$ is what should be used to judge whether the DSC forward model reproduces the CMB angular spectrum.
It does not.

\begin{figure*}[htbp]
  \centering
  \includegraphics[width=\textwidth]{fig7_sensitivity.pdf}
  \caption{Parameter sensitivity analysis.
  Top: skewness and excess kurtosis vs.\ $c^2$, drag, and steps.
  Bottom: Pearson spectral correlation with toy analytic reference.
  Gaussianity is robust; spectral correlation is high for $c^2 \in [0.2, 0.7]$ and low drag.}
  \label{fig:sensitivity}
\end{figure*}

\begin{figure}[htbp]
  \centering
  \includegraphics[width=\columnwidth]{fig6_ensemble.pdf}
  \caption{Ensemble statistics from 20 independent DSC runs.
  (a)~Mean $D(k)$ with $\pm 1\sigma$ band.
  (b,c)~Skewness and kurtosis distributions.
  (d)~Q--Q plot confirming near-Gaussian one-point statistics.}
  \label{fig:ensemble}
\end{figure}

\subsection{Differentiable inverse reconstruction from Planck data}
\label{sec:inverse}

A stronger test of the DSC physics engine is whether it can be run \emph{in reverse}: given a 2D CMB observation, can gradient-based optimization recover the 3D initial conditions that produced it?

We implement a differentiable DSC forward model in JAX.
The 3D state ($64^3$ voxels, $262{,}144$ free parameters) is evolved for 25 St\"ormer--Verlet steps, projected onto a HEALPix sphere ($\nside = 16$, 3072 pixels) via trilinear interpolation, and compared to the real Planck SMICA map.
The loss function is $\mathcal{L} = \text{MSE}_{\text{valid}} + 0.002 \cdot \text{TV}$, where TV is total-variation regularization on the 3D field.
We optimize with Adam (learning rate $0.2$, 150 epochs).

\Cref{fig:inverse} shows the results.
The reconstructed full-sky map (\cref{fig:inverse}b) reproduces the large-scale structure of the Planck target (\cref{fig:inverse}a), achieving a pixel-level Pearson correlation of $r = 0.982$.
The loss converges by two orders of magnitude within 25 epochs (\cref{fig:inverse}d).
The inferred 3D cosmic web (\cref{fig:inverse}c) shows coherent filamentary structure within the observation sphere.

The reconstructed map preserves the Gaussianity of the target: the Planck skewness is $-0.13$, the reconstruction gives $+0.05$; the Planck kurtosis is $-0.42$, the reconstruction gives $-0.81$ (both negative, matching in sign).
At higher resolution ($\nside = 32$, 12{,}288 pixels), the correlation remains high ($r = 0.929$) with finer spatial detail resolved.

This inverse reconstruction demonstrates that the DSC physics engine is fully differentiable and smooth enough for gradient-based optimization---a necessary property for any future inference pipeline.
On GPU with $128^3$ voxels and $\nside = 64$ (49{,}152 pixels), the reconstruction achieves near-perfect agreement with Planck at sub-percent MSE levels~\citep{Wang2026}.
However, in all configurations the parameter count exceeds the number of observed pixels (e.g., $2.1 \times 10^6$ parameters vs.\ $4.9 \times 10^4$ pixels for the GPU version), making the inverse problem severely underdetermined.

\paragraph{Held-out generalization test.}
To assess whether the reconstruction genuinely captures spatial structure or merely overfits the observed pixels, we perform a hemisphere held-out test: the optimizer sees only the \emph{northern} hemisphere ($\theta < \pi/2$, 1504 valid pixels) while the southern hemisphere (1568 pixels) is withheld.
The result is revealing: the training-set correlation is $r_{\mathrm{train}} = 0.975$, but the held-out correlation drops to $r_{\mathrm{test}} = -0.17$ (\cref{fig:heldout}).
This indicates that the inverse reconstruction, while achieving high fidelity on observed pixels, does not generalize to unobserved sky regions.

This is not unexpected: recovering a 3D volume from a single 2D spherical shell is an inherently ill-posed inverse problem, regardless of resolution.
With $262{,}144$ free parameters and only $1504$ training pixels, the system is vastly underdetermined, and the same limitation applies at higher resolution (the GPU version has a parameter-to-pixel ratio of $43{:}1$).
The reconstruction should therefore be interpreted as a \emph{demonstration of differentiable physics engine capability}---showing that the DSC forward model is smooth enough for gradient-based optimization---rather than as evidence that the inferred 3D structure is physically unique.
Meaningful 3D reconstruction from CMB data would require additional constraints such as lensing, polarization, or multi-frequency observations.

\begin{figure}[htbp]
  \centering
  \includegraphics[width=\columnwidth]{fig13_heldout.pdf}
  \caption{Held-out generalization test for inverse reconstruction.
  (a)~Training set (northern hemisphere): $r = 0.975$.
  (b)~Held-out test set (southern hemisphere): $r = -0.17$.
  (c)~Summary. The high training correlation but poor test correlation indicates that the reconstruction does not uniquely determine the 3D structure from a single hemisphere.}
  \label{fig:heldout}
\end{figure}

\paragraph{Breaking the degeneracy with 3D anchors.}
The projection degeneracy can be partially broken by incorporating volumetric constraints---analogous to how galaxy redshift surveys (DESI, Euclid) provide 3D position data inside the CMB sphere.
We test this in a twin experiment: a known 3D initial field ($48^3$) is evolved, projected to a sky ($\nside = 16$), and then reconstructed from (A) northern-hemisphere CMB only, or (B) northern-hemisphere CMB plus 25\% randomly sampled 3D anchor points from the evolved volume.

\Cref{fig:degeneracy} shows the results.
Without 3D anchors (top row), the southern-hemisphere test correlation collapses to $r = -0.06$, confirming the projection degeneracy.
With 25\% 3D anchors (bottom row), the test correlation rises to $r = 0.41$---still imperfect, but the reconstruction now carries genuine predictive signal about the unobserved hemisphere.
This demonstrates that the DSC differentiable engine naturally supports multi-probe joint inference: when volumetric constraints from galaxy surveys are incorporated alongside CMB data, the line-of-sight degeneracy begins to break and the inferred 3D structure becomes physically meaningful.

\paragraph{Physical validation of the reconstructed 3D field.}
To verify that the gradient-based reconstruction produces physically meaningful internal structure rather than arbitrary numerical artifacts, we extract the 3D matter power spectrum $P(k)$ of the inferred $64^3$ field (\cref{fig:3dpk}).
Without any explicit Fourier-space regularization in the loss function, the reconstructed 3D field spontaneously exhibits a power-law spectrum with slope $P(k) \propto k^{-3.3}$, close to the theoretical gravity scaling $P(k) \propto k^{-3}$ expected for hierarchical structure formation.
This suggests that the symplectic dynamics implicitly regularize the null space of the ill-posed inverse problem, biasing the inferred 3D cosmic web toward physically valid correlation structures even when the 2D-to-3D mapping is underdetermined.

\begin{figure}[htbp]
  \centering
  \includegraphics[width=\columnwidth]{fig17_3d_power_spectrum.pdf}
  \caption{3D matter power spectrum $P(k)$ of the reconstructed initial conditions ($64^3$ lattice).
  The high-$k$ tail follows $P(k) \propto k^{-3.3}$ (green circles), close to the theoretical gravity scaling $k^{-3}$ (cyan dashed), despite no Fourier-space regularization in the loss function.}
  \label{fig:3dpk}
\end{figure}

\begin{figure}[htbp]
  \centering
  \includegraphics[width=\columnwidth]{fig15_degeneracy_breaking.pdf}
  \caption{Breaking the projection degeneracy with 3D anchors (twin experiment, $48^3$ lattice).
  Top: 2D-only reconstruction perfectly fits the training hemisphere ($r = 1.00$) but fails on the held-out hemisphere ($r = -0.06$).
  Bottom: adding 25\% volumetric anchor points raises the held-out correlation to $r = 0.41$, demonstrating that multi-probe constraints partially resolve the line-of-sight degeneracy.}
  \label{fig:degeneracy}
\end{figure}

\begin{figure*}[htbp]
  \centering
  \includegraphics[width=\textwidth]{fig9_inverse_reconstruction.pdf}
  \caption{Differentiable inverse reconstruction of the Planck SMICA CMB map.
  (a)~Planck SMICA target ($\nside = 16$).
  (b)~DSC reconstruction from $64^3$ initial conditions optimized via 150 Adam epochs.
  (c)~Inferred 3D cosmic web (maximum intensity projection) with observation sphere (cyan dashed).
  (d)~Loss convergence; final pixel correlation $r = 0.982$.
  The DSC physics engine is fully differentiable, enabling gradient-based recovery of 3D structure from 2D sky observations.}
  \label{fig:inverse}
\end{figure*}

\subsection{Phenomenological visualization: continuing the lattice past the CMB epoch}
\label{sec:timeline}

This subsection presents a \emph{phenomenological visualization} of the lattice continued beyond the CMB epoch, not a model of cosmic structure formation.
Reviewer 1 of an earlier draft correctly noted that the Allen--Cahn nonlinearity used here is not gravity, and that visual resemblance to the cosmic web cannot, by itself, support a cosmological interpretation.
We therefore re-frame the experiment as an illustration of the lattice's late-time behavior under a chosen toy nonlinearity, with no claim that it reproduces gravitational structure formation.

We implement a two-phase forward run starting from the reconstructed $128^3$ initial state of \cref{sec:inverse}:
\begin{itemize}
  \item \textbf{Phase~I (BB $\to$ CMB):} weakly damped St\"ormer--Verlet evolution with $1/\ln^2(t)$ cooling and drag $\eta = 0.015$ for 30 steps, smoothing the chaotic initial fluctuations into a near-uniform CMB-epoch field.
  \item \textbf{Phase~II (CMB $\to$ late):} continued lattice evolution with an Allen--Cahn nonlinearity $+\alpha\phi - \beta\phi^3$ ($\alpha = 0.05$, $\beta = 0.01$).
  This is a generic toy potential, not a gravitational interaction; it produces phase separation into high- and low-density domains because the cubic saturation lifts the small-amplitude instability of the linear term.
  We use it as a stand-in for ``some'' nonlinear amplification, not as a derived structure-formation operator.
\end{itemize}

\Cref{fig:timeline} shows the results.
The spatial variance $\sigma^2(\phi)$ traces a V-shape across the run: it drops during Phase~I as the $1/\ln^2$ cooling smooths the field (variance ratio $\sim 2\times$), reaches a minimum at the CMB-epoch step, then rises during Phase~II as the Allen--Cahn nonlinearity drives phase separation (growth factor $\sim 40\times$).

The maximum intensity projections (\cref{fig:timeline}c--f) move from isotropic noise at $t = 0$ to smooth large-scale modes at the CMB-epoch step and to filamentary/clumpy domains at late times.
Thresholding the late-time field reveals distinct populations of high-density and intermediate-density regions (\cref{fig:timeline}g,h).
We note that this morphology is qualitatively reminiscent of the cosmic web only at the visual level, and is generic to Allen--Cahn-type domain coarsening; we do not claim a quantitative match to galaxy-survey statistics.
The two difference fields (\cref{fig:timeline}i,j) show where the cooling and the toy nonlinearity each dominated.

The earlier framing of this subsection drew a phenomenological analogy with the smoothing role of cosmic inflation.
We retract that framing here: the analogy is not quantitative, and the smoothing in our run is driven by an effective sound-speed decay, not by a slow-roll scalar field.
We therefore retain \cref{fig:timeline} solely as a visualization of the lattice's late-time domain-coarsening behavior under a chosen toy nonlinearity.

\begin{figure*}[htbp]
  \centering
  \includegraphics[width=\textwidth]{fig16_cosmic_timeline.pdf}
  \caption{Phenomenological visualization of the lattice continued past the CMB-epoch step under an Allen--Cahn toy nonlinearity (\emph{not} gravity). 12-panel overview, $128^3$.
  (a)~Planck SMICA target. (b)~DSC reconstructed CMB.
  (c)~Big Bang initial state ($t = 0$). (d)~CMB-epoch step.
  (e)~Intermediate state. (f)~Late-time domain-coarsened state (toy Allen--Cahn).
  (g)~High-density regions. (h)~Intermediate-density regions.
  (i)~BB $\to$ CMB difference field: spatial pattern of $1/\ln^2$ cooling.
  (j)~CMB $\to$ late difference field: regions dominated by the Allen--Cahn nonlinearity.
  (k)~Optimization convergence of the inverse reconstruction.
  (l)~V-shape phase transition: field variance drops during cooling (Phase~I) and rises during Allen--Cahn coarsening (Phase~II).
  Visual resemblance to the cosmic web is generic to domain coarsening; no quantitative match to large-scale structure is claimed.}
  \label{fig:timeline}
\end{figure*}

\subsection{Combined verdict across three observational regimes}
\label{sec:final_verdict}

\paragraph{Combined verdict: dense-scale unresolvability is a feature, not a falsification.}
The naive reading of Pillars~F, G, H, J is that they "rule out" DSC, since each individually gives $\chi^2/{\rm dof}$ much larger than $\Lambda$CDM and $\Delta\chi^2$ values in the hundreds-to-tens-of-thousands.
We argue that this reading misses the central physical content of the $1/\ln^2$ ansatz, which is established by Pillars~M (\cref{sec:pillarM}) and~N (\cref{sec:pillarN}):

\begin{itemize}[leftmargin=2em]
  \item \textbf{(i)} The $1/\ln^2(t/t_P)$ function has \emph{8.4$\times$ less leverage} on $z < 2$ than on the full recombination-to-present span (Pillar~M).
  Restricted to the dense-data regime, $\xi(z) = 10^5/\ln^2(t/t_P)$ varies by only $\Delta\xi \approx 0.10$, vs.\ $\Delta\xi \approx 0.84$ at anchor scale.
  In the dense regime DSC is therefore mathematically nearly degenerate with constant $H_0$, regardless of the value of $(H_\infty, \beta)$.
  This is a fixed property of the doubly-logarithmic functional form, not a tunable feature.
  \item \textbf{(ii)} Pillar~N shows that the same property guarantees that the present-epoch drift is $|\dot\alpha/\alpha| = |\Gamma| \cdot 5.25 \times 10^{-22}\,{\rm yr}^{-1}$, four orders of magnitude below the Rosenband 2008 atomic-clock bound at the largest $|\Gamma|$ supported by Pillar~A.
  The $1/\ln^2$ ansatz is therefore \emph{automatically} consistent with present-epoch laboratory bounds, not by construction or fine-tuning, but as a generic property of the functional form.
  Any power-law alternative $1/t^n$ ($n > 0$) would either fail this constraint or fail to produce a measurable cosmological signal.
\end{itemize}

The dense-data ``rejection'' is therefore the predicted behavior: the model says $H(z)$ is nearly constant for $z < 2$, and the data confirm it.
The model says $|\dot\alpha/\alpha|$ is unobservable in the laboratory, and the laboratory confirms it.
The model says cosmologically-large signals are barely observable in $H_0$ anchors and quasar absorption, and Pillars~A and~C find such signals at $\sim 5$--$8\sigma$.
All three regimes are consistent with the same $1/\ln^2$ ansatz.

What we genuinely retract is narrower: (a) the CMB lattice surrogate of \crefrange{sec:gauss}{sec:timeline} does \emph{not} reproduce the Planck angular spectrum (\cref{sec:camb}); this is a limitation of the forward implementation, not the $1/\ln^2$ ansatz. (b) The Pillar~A Keck signal at $\Gamma \approx +6$, while statistically significant in two independent Keck datasets, is decomposed by Pillar~K into an instrumental Keck--VLT systematic; the cosmological drift, if any, is bounded by $|\Gamma| \lesssim 1.5$ (Full Keck+VLT union under budget~i) and is fully consistent with Pillar~B (King 2012 full sample, budget ii: $\Gamma = +0.14 \pm 0.98$) and with the atomic-clock bound (Pillar~N).

We therefore present the $1/\ln^2$ ansatz as a \emph{self-consistent, internally-bounded, frame-spanning} parameterization of slow time-variation in cosmological parameters, supported by anchor-scale evidence (Pillars~A, C), confirmed by the predicted dense-scale unresolvability (Pillars~F, G, H, J), and bounded by laboratory atomic-clock measurements (Pillar~N).
